\section{System in Use}
\label{sec:demo}

Le Rapporteur is deployed as a working, public demonstration, deliberately
styled after the French State Design System (\emph{Syst\`eme de design de
l'\'Etat})---sober typography, the tricolour, no chatbot theatrics---so that the
interface itself signals institutional trust rather than novelty. Three views
structure the experience: \emph{ask} (pose a question and receive a sourced
answer or a refusal), \emph{how it works} (the pipeline of
\S\ref{sec:overview}, made legible to a non-specialist), and \emph{sources}
(browse the verbatim, in-force article text behind an answer).

\paragraph{What a visitor sees.}
For an answerable question, the answer appears with each citation rendered as a
clickable link to the exact \texttt{LEGIARTI} version on L\'egifrance, next to
the excerpt the verifier returned---provenance a citizen can check in a single
click. For a trap question, the interface shows the explicit refusal in place
of a fabricated article, making the fail-closed behaviour \emph{visible} rather
than hidden in a log. A frugality panel foregrounds the
sovereignty-and-efficiency argument of \S\ref{sec:sovereign}, so that
perf/watt is part of the story a decision-maker sees, not a footnote.

\paragraph{Faithful static builds.}
For zero-backend hosting, the same pipeline precomputes answers for a fixed
question set (\S\ref{sec:service}), so the public demonstration is generated by
the real system rather than re-implemented by hand. The trust
story---sourced answers, red/green refusals, clickable primary sources---is
therefore identical whether the demo is served live from the FastAPI backend or
statically from precomputed artifacts.
